% steps/step12.tex
\section[گام دوازدهم: پایش و مانیتورینگ]{\faChartBar\ \ گام دوازدهم: پایش و مانیتورینگ}
\begin{stepbox}

\field{پایش عملکرد در دنیای واقعی (Model Monitoring)}{
    پس از استقرار سیستم، ارزیابی مستمر دقت مدل در مواجهه با بیماران جدید با ویژگی‌های بالینی و سنی متفاوت ضروری است. این کار از طریق پایش غیرفعال و دریافت بازخوردهای اصلاحی مستقیم از رادیولوژیست‌ها (به‌عنوان مثال، ثبت میزان اصلاح دستی ماسک پیشنهادی مدل توسط پزشک در سامانه \lr{PACS}) انجام می‌شود تا هرگونه افت کارایی مدل در بلندمدت شناسایی شود.
}

\field{بررسی مداوم کیفیت داده‌های ورودی (Data Drift \& Shift)}{
    تغییر سخت‌افزار اسکنرها، اعمال پروتکل‌های جدید تصویربرداری یا تغییر غلظت ماده حاجب، توزیع داده‌های ورودی را نسبت به داده‌های زمان آموزش تغییر می‌دهد (\lr{Data Drift}). سیستم مانیتورینگ باید به‌صورت خودکار شاخص‌های آماری تصاویر جدید را رصد کرده و در صورت بروز انحراف جدی از الگوی داده‌های اولیه، به تیم فنی هشدار دهد تا از تحلیل نادرست تصاویر جلوگیری شود.
}

\field{ثبت و تحلیل خطاها و هشدارها (Logging \& Alerting)}{
    تمام تراکنش‌ها، فراداده‌های تصاویر ورودی، نتایج پیش‌بینی و زمان پردازش مدل باید در بستری امن ثبت (\lr{Log}) شوند. در صورت بروز هرگونه مشکل سیستمی مانند افزایش زمان استنتاج، خطاهای خارج از محدوده حافظه پردازنده گرافیکی یا توقف ناگهانی سرویس، مکانیزم‌های هشداردهی باید بلافاصله فعال شده و تیم پشتیبانی فنی را مطلع سازند.
}

\end{stepbox}
